\section{Evaluation}
\label{sec:evaluation}

We evaluate multi-member districts on four dimensions across three case studies: California (52 seats), New Jersey (12 seats), and Illinois (17 seats), all using 2020 presidential election data.

\subsection{RQ1: Gerrymandering Elimination}

Table~\ref{tab:mmid-efficiency-gap} reports the efficiency gap under three district configurations for each state.

\begin{table}[h]
\centering\small
\caption{Efficiency gap by district configuration. Lower absolute values indicate better partisan fairness. All single-member results use standard recursive bisection (geographic edge weights, no partisan data).}
\label{tab:mmid-efficiency-gap}
\begin{tabular}{lcccc}
\toprule
State & Single-Member (enacted) & Single-Member (algo) & 4-Member MMD & 3-Member MMD \\
\midrule
California  & $+8.1\%$ (R) & $-3.1\%$ (D) & $-0.6\%$ & $-0.5\%$ \\
New Jersey  & $+2.3\%$ (R) & $-2.4\%$ (D) & $-0.3\%$ & $-0.2\%$ \\
Illinois    & $-5.1\%$ (D) & $-2.9\%$ (D) & $-0.4\%$ & $-0.4\%$ \\
\bottomrule
\end{tabular}
\end{table}

MMDs reduce the efficiency gap to below $\pm 0.7\%$ in all three states — below the measurement precision of precinct-level voting data. The residual Democratic advantage reflects geographic sorting (urban Democratic concentration) rather than algorithmic bias: single-member algorithmic plans also show small Democratic efficiency gaps from the same geographic sorting effect. MMDs suppress rather than eliminate this bias because proportional seat allocation within each district forces representation of the minority party even in partisan-homogeneous districts.

\subsection{RQ2: Minority Representation}

Table~\ref{tab:mmid-minority} compares majority-minority districts under single-member and MMD configurations.

\begin{table}[h]
\centering\small
\caption{Majority-minority seats by configuration. MMD counts are ``effective'' minority opportunity seats where minority voters $\geq 40\%$ within at least one district and constitute $> 1$ seat equivalent of representation.}
\label{tab:mmid-minority}
\begin{tabular}{lcccc}
\toprule
State & Enacted MM & Algo MM & 4-Member MMD & 3-Member MMD \\
\midrule
California  & 17 & 19 & 22 & 21 \\
New Jersey  &  2 &  3 &  4 &  3 \\
Illinois    &  4 &  5 &  7 &  6 \\
\bottomrule
\end{tabular}
\end{table}

MMDs improve minority opportunity seat counts by 15--30\% over algorithmic single-member plans and 25--50\% over enacted maps. The improvement arises from fractional representation: a district with 35\% minority population elects approximately 1.4 of 4 MMD seats for the minority-preferred candidate under OLPR rules, compared to 0 seats under winner-take-all single-member districts.

\subsection{RQ3: Optimal District Size}

Figure~\ref{fig:district-size-tradeoff} shows the Pareto frontier between compactness and proportionality as the district size $m$ varies from 2 to 7 members.

The key tradeoff: larger districts improve proportionality (higher effective thresholds tolerate more diverse within-district vote shares) but reduce compactness (each district must span a larger geographic area). From the data:

\begin{itemize}
  \item $m=2$ (two-member): modest proportionality gain, only 4\% compactness loss
  \item $m=3$: preferred by most states — gains 80\% of the proportionality benefit of $m=7$ at only 8\% compactness loss
  \item $m=4$: used in California case study — good balance, 10\% compactness loss
  \item $m=7$: maximum proportionality, 18\% compactness loss, states become city-sized
\end{itemize}

Based on this analysis, 3- to 4-member districts are the recommended configuration for states seeking improved proportionality while preserving geographic compactness to a degree defensible under constitutional standards.

\subsection{RQ4: Political Competition and Third-Party Viability}

The natural threshold for third-party representation is $1/(m+1)$: a party needs $> 25\%$ of votes within a district to win one of four seats in a four-member district. Table~\ref{tab:mmid-threshold} shows thresholds by district size.

\begin{table}[h]
\centering\small
\caption{Natural threshold for third-party seat by district size. Below the threshold, a third party cannot win any seat in that district even with maximum vote efficiency.}
\label{tab:mmid-threshold}
\begin{tabular}{lcccccc}
\toprule
$m$ (members) & 2 & 3 & 4 & 5 & 6 & 7 \\
\midrule
Natural threshold & 33.3\% & 25.0\% & 20.0\% & 16.7\% & 14.3\% & 12.5\% \\
\bottomrule
\end{tabular}
\end{table}

At $m=4$, a third party winning 20\% of votes in a district wins one of four seats in that district. California's 13 four-member districts at 20\% threshold would yield 3--4 third-party seats if a third party won 20\% statewide — compared to zero under the enacted winner-take-all system where third parties win 8\% of the statewide two-party-equivalent vote (2020).

Cross-census stability (2000/2010/2020) shows that MMD district configurations are more stable than single-member maps across census cycles: mean RSI across the three censuses is 0.04 for four-member California districts vs.\ 0.12 for single-member California districts, reflecting the fact that larger districts absorb population shifts without requiring boundary changes.
